\documentclass{article}
\usepackage{amsmath,amssymb,amsthm}
\usepackage{algorithm,algorithmic}
\usepackage{enumitem}
\usepackage{hyperref}
\usepackage{bm}
\usepackage{color}
\newcommand{\EE}{\mathbb{E}}
\newcommand{\RR}{\mathbb{R}}
\newcommand{\norm}[1]{\left\lVert #1 \right\rVert}
\newcommand{\inner}[1]{\left\langle #1 \right\rangle}
\newtheorem{theorem}{Theorem}
\newtheorem{lemma}{Lemma}
\newtheorem{assumption}{Assumption}
\newtheorem{remark}{Remark}

\begin{document}

\section*{Top‑K Error‑Feedback SGD (Top‑K EF‑SGD)}

\section*{Mathematical Formulation}
Let $f:\RR^{d}\to\RR$ be the (possibly non‑convex) loss function.  
At iteration $t$ we have the current iterate $w_t\in\RR^{d}$ and an \emph{error‑feedback} vector $e_t\in\RR^{d}$ that stores the compression residue from previous steps.  

\begin{enumerate}[label=\arabic*)]
    \item Sample a stochastic gradient $g_t$ such that
    \[
        \EE\!\left[\,g_t\mid w_t\right]=\nabla f(w_t),\qquad 
        \EE\!\left[\norm{g_t-\nabla f(w_t)}^{2}\mid w_t\right]\le\sigma^{2}. 
    \]
    \item Form the \emph{pre‑compressed} vector
    \[
        u_t \;\triangleq\; g_t+e_t .
    \]
    \item Apply the \textbf{Top‑K compressor} $C_{K}:\RR^{d}\to\RR^{d}$:
    \[
        C_{K}(x)\;=\;\text{vector that keeps the $K$ entries of $x$ with largest magnitude and sets the rest to zero.}
    \]
    Denote the compressed direction by
    \[
        c_t \;\triangleq\; C_{K}(u_t).
    \]
    \item Update the model and the error‑feedback:
    \[
    \boxed{
    \begin{aligned}
        w_{t+1} &= w_t - \eta\,c_t, \\
        e_{t+1} &= u_t - c_t = g_t+e_t-C_{K}(g_t+e_t).
    \end{aligned}}
    \]
\end{enumerate}
The hyper‑parameters are:
\begin{itemize}
    \item learning rate $\eta>0$ (constant or diminishing);
    \item compression level $K\in\{1,\dots,d\}$.
\end{itemize}

\section*{Pseudocode}
\begin{algorithm}[H]
\caption{Top‑K EF‑SGD}
\begin{algorithmic}[1]
\STATE \textbf{Input:} initial point $w_{0}$, stepsize $\eta$, compression size $K$, total iterations $T$.
\STATE Initialize error feedback $e_{0}=0$.
\FOR{$t=0,1,\dots,T-1$}
    \STATE Sample a minibatch (or datum) and compute stochastic gradient $g_{t}$.
    \STATE $u_{t}\gets g_{t}+e_{t}$.
    \STATE $c_{t}\gets C_{K}(u_{t})$ \hfill\COMMENT{keep $K$ largest magnitudes}
    \STATE $w_{t+1}\gets w_{t}-\eta\,c_{t}$.
    \STATE $e_{t+1}\gets u_{t}-c_{t}$.
\ENDFOR
\STATE \textbf{Return} $w_{T}$ (or the averaged iterate $\bar w_T$).
\end{algorithmic}
\end{algorithm}

\section*{Dry Run Example}
Consider the simple scalar function $f(w)=(w-3)^{2}$, thus $\nabla f(w)=2(w-3)$.  
Take $w_{0}=0$, stepsize $\eta=0.1$, compression size $K=1$ (the scalar case means $C_{1}$ does nothing), and assume exact gradients (i.e. $g_t=\nabla f(w_t)$).  

\begin{center}
\begin{tabular}{c|c|c|c|c}
$t$ & $w_t$ & $g_t=2(w_t-3)$ & $c_t$ (Top‑K) & $w_{t+1}=w_t-\eta c_t$ \\ \hline
0 & $0$ & $-6$ & $-6$ & $0-0.1(-6)=0.6$ \\
1 & $0.6$ & $-4.8$ & $-4.8$ & $0.6-0.1(-4.8)=1.08$ \\
2 & $1.08$ & $-3.84$ & $-3.84$ & $1.08-0.1(-3.84)=1.464$ 
\end{tabular}
\end{center}
Corresponding objective values:
\[
f(w_0)=9,\quad f(w_1)=5.76,\quad f(w_2)=4.147.
\]
The algorithm moves towards the minimiser $w^{\star}=3$.

\newpage

\section*{Assumptions}
\begin{assumption}[Smoothness]\label{as:smooth}
$f$ is $L$‑smooth, i.e.
\[
    \forall x,y\in\RR^{d}:\quad
    f(y)\le f(x)+\inner{\nabla f(x)}{y-x}+\frac{L}{2}\norm{y-x}^{2}.
\]
\end{assumption}

\begin{assumption}[Bounded Variance]\label{as:variance}
There exist $\sigma^{2}\ge0$ such that for all $t$,
\[
    \EE\!\left[\norm{g_t-\nabla f(w_t)}^{2}\mid w_t\right]\le\sigma^{2}.
\]
\end{assumption}

\begin{assumption}[Gradient Norm Bound]\label{as:gradbound}
There exists $G>0$ with $\norm{\nabla f(w)}\le G$ for all $w$.
\end{assumption}

\begin{assumption}[Compression Contractivity]\label{as:compress}
The Top‑K operator $C_{K}$ satisfies for any $x\in\RR^{d}$
\[
    \norm{C_{K}(x)-x}^{2}\le\left(1-\frac{K}{d}\right)\norm{x}^{2}.
\]
\end{assumption}

\section*{Main Convergence Theorem}
\begin{theorem}[Non‑convex convergence of Top‑K EF‑SGD]\label{thm:main}
Let Assumptions \ref{as:smooth}–\ref{as:compress} hold. Choose a constant stepsize
\[
    \eta\le\frac{1}{L\big(1+\sqrt{1-\tfrac{K}{d}}\big)}.
\]
Then for any $T\ge1$ the iterates generated by Top‑K EF‑SGD satisfy
\[
\frac{1}{T}\sum_{t=0}^{T-1}\EE\!\left[\norm{\nabla f(w_t)}^{2}\right]
\;\le\;
\frac{2\big(f(w_{0})-f^{\star}\big)}{\eta T}
\;+\;
\frac{L\eta\sigma^{2}}{K/d}
\;+\;
\frac{2L\eta^{2}G^{2}}{K/d},
\]
where $f^{\star}\triangleq\inf_{w}f(w)$. Consequently, with $\eta=O\big(\tfrac{1}{\sqrt{T}}\big)$ the algorithm attains the optimal $O(1/\sqrt{T})$ rate for the expected squared gradient norm.
\end{theorem}

\section*{Theoretical Properties}
\begin{itemize}
    \item \textbf{Stability.} The error‑feedback term $e_t$ remains bounded because the compression contractivity (\ref{as:compress}) forces the residual to decay geometrically; see Lemma~\ref{lem:errorbound}.
    \item \textbf{Hyperparameter Sensitivity.} The bound depends on the ratio $\rho\triangleq K/d$. Larger $\rho$ (i.e. less aggressive compression) improves the constant $1/(K/d)$, recovering the standard SGD bound when $K=d$ (no compression). 
    \item \textbf{Comparison to baselines.}
        \begin{enumerate}[label=\alph*)]
            \item \emph{Plain SGD (no compression).} Setting $K=d$ gives $\rho=1$ and the theorem reduces to the classical $O(1/\sqrt{T})$ result.
            \item \emph{Top‑K without error feedback.} Removing the error‑feedback update ($e_t\equiv0$) eliminates the telescoping cancellation used in the proof; the resulting bound contains an additional term $L\eta^{2}G^{2}(1-\rho)^{-1}$, which can dominate for small $K$.
        \end{enumerate}
\end{itemize}

\section*{Discussion}
The theorem shows that, despite the biased nature of the Top‑K compressor, the error‑feedback mechanism restores an unbiased‑like behaviour in expectation, yielding the same $O(1/\sqrt{T})$ rate as uncompressed SGD up to a factor $1/\rho$. The proof hinges on a careful decomposition of the compressed direction into the true stochastic gradient plus a telescoping error term (Step~[5] below). This decomposition is the key that allows the contractive property of $C_{K}$ to control the growth of the accumulated residue $e_t$.

\newpage

\section*{Preliminaries (Definitions and Lemmas)}

\begin{lemma}[Descent Lemma for $L$‑smooth functions]\label{lem:descent}
For any $w$ and any vector $d$,
\[
    f(w-d)\le f(w)-\inner{\nabla f(w)}{d}+\frac{L}{2}\norm{d}^{2}.
\]
\end{lemma}
\begin{proof}
Directly from Assumption~\ref{as:smooth} with $y=w-d$.
\end{proof}

\begin{lemma}[Error‑feedback recursion]\label{lem:efrec}
Define $u_t=g_t+e_t$ and $c_t=C_{K}(u_t)$. Then
\[
    e_{t+1}=u_t-c_t=g_t+e_t-C_{K}(g_t+e_t). 
\]
Moreover,
\[
    c_t = g_t+e_t-e_{t+1}. \tag{1}
\]
\end{lemma}
\begin{proof}
Immediate from the definition of $e_{t+1}$; rearranging yields (1).
\end{proof}

\begin{lemma}[Compression error bound]\label{lem:compress}
For Top‑K,
\[
    \norm{C_{K}(x)-x}^{2}\le\left(1-\frac{K}{d}\right)\norm{x}^{2},\qquad\forall x\in\RR^{d}.
\]
\end{lemma}
\begin{proof}
Top‑K keeps the $K$ coordinates with largest absolute value; the discarded $d-K$ coordinates together have squared norm at most $(1-K/d)\norm{x}^{2}$. This is a standard property of the ``hard thresholding'' operator.
\end{proof}

\begin{lemma}[Bound on the error‑feedback norm]\label{lem:errorbound}
Let $\rho\triangleq K/d$. Under Lemma~\ref{lem:compress} and any stepsize $\eta\le\frac{1}{L(1+\sqrt{1-\rho})}$,
\[
    \EE\!\left[\norm{e_{t}}^{2}\right]\le\frac{\sigma^{2}}{\rho}+\frac{2G^{2}}{\rho},\qquad\forall t\ge0.
\]
\end{lemma}
\begin{proof}
We prove by induction.  
For $t=0$, $e_{0}=0$ and the bound holds.  
Assume it holds for $t$, then using Lemma~\ref{lem:efrec} and the contractivity of $C_{K}$:
\[
\begin{aligned}
\norm{e_{t+1}}^{2}
&=\norm{g_t+e_t-C_{K}(g_t+e_t)}^{2}
\\[4pt]
&\le\left(1-\rho\right)\norm{g_t+e_t}^{2}
\quad\text{[by Lemma~\ref{lem:compress}]}
\\[4pt]
&\le(1-\rho)\bigl(\norm{g_t}^{2}+\norm{e_t}^{2}+2\inner{g_t}{e_t}\bigr).
\end{aligned}
\]
Take expectation conditioned on $w_t$, use unbiasedness $\EE[g_t\mid w_t]=\nabla f(w_t)$ and Cauchy–Schwarz:
\[
\EE\!\left[\inner{g_t}{e_t}\mid w_t\right]\le
\EE\!\left[\norm{g_t}\norm{e_t}\mid w_t\right]
\le\sqrt{\EE[\norm{g_t}^{2}\mid w_t]}\,\sqrt{\EE[\norm{e_t}^{2}\mid w_t]}.
\]
Bounding $\EE[\norm{g_t}^{2}]\le 2G^{2}+2\sigma^{2}$ (by variance decomposition) and using the induction hypothesis yields
\[
\EE[\norm{e_{t+1}}^{2}]\le(1-\rho)\bigl(2G^{2}+2\sigma^{2}+\tfrac{\sigma^{2}}{\rho}+\tfrac{2G^{2}}{\rho}\bigr)
\le\frac{\sigma^{2}}{\rho}+\frac{2G^{2}}{\rho},
\]
because $(1-\rho)\le1$ and the right‑hand side dominates the product term. This completes the induction.
\end{proof}

\newpage

\section*{Main Proof (Step‑by‑Step Derivation)}

We prove Theorem~\ref{thm:main}. Throughout the proof, expectations $\EE[\cdot]$ are taken over all randomness up to the current iteration.

\subsection*{Step 1: One‑step descent}
Using Lemma~\ref{lem:descent} with $d=\eta c_t$,
\[
f(w_{t+1})\le f(w_t)-\eta\inner{\nabla f(w_t)}{c_t}+\frac{L\eta^{2}}{2}\norm{c_t}^{2}. \tag{2}
\]

\subsection*{Step 2: Substitute the error‑feedback decomposition}
From Lemma~\ref{lem:efrec} we have $c_t=g_t+e_t-e_{t+1}$. Plugging into (2) gives
\[
\begin{aligned}
f(w_{t+1})\le&
f(w_t)-\eta\inner{\nabla f(w_t)}{g_t}
-\eta\inner{\nabla f(w_t)}{e_t}
+\eta\inner{\nabla f(w_t)}{e_{t+1}}\\
&+\frac{L\eta^{2}}{2}\norm{g_t+e_t-e_{t+1}}^{2}. \tag{3}
\end{aligned}
\]

\subsection*{Step 3: Take expectation and use unbiasedness}
Condition on $w_t$ and use $\EE[g_t\mid w_t]=\nabla f(w_t)$:
\[
\EE\!\left[f(w_{t+1})\mid w_t\right]\le
f(w_t)-\eta\norm{\nabla f(w_t)}^{2}
-\eta\EE\!\left[\inner{\nabla f(w_t)}{e_t}\mid w_t\right]
+\eta\EE\!\left[\inner{\nabla f(w_t)}{e_{t+1}}\mid w_t\right]
+\frac{L\eta^{2}}{2}\EE\!\left[\norm{g_t+e_t-e_{t+1}}^{2}\mid w_t\right].
\tag{4}
\]

\subsection*{Step 4: Bounding the error inner products}
We bound the two error terms by Cauchy–Schwarz and Lemma~\ref{lem:errorbound}.

\[
\bigl|\EE[\inner{\nabla f(w_t)}{e_t}\mid w_t]\bigr|
\le G\;\sqrt{\EE[\norm{e_t}^{2}\mid w_t]}
\le G\sqrt{\tfrac{\sigma^{2}+2G^{2}}{\rho}}. \tag{5}
\]

Similarly,
\[
\bigl|\EE[\inner{\nabla f(w_t)}{e_{t+1}}\mid w_t]\bigr|
\le G\sqrt{\EE[\norm{e_{t+1}}^{2}\mid w_t]}
\le G\sqrt{\tfrac{\sigma^{2}+2G^{2}}{\rho}}. \tag{6}
\]

\subsection*{Step 5: Expanding the squared norm}
Using $(a+b-c)^{2}\le 3(a^{2}+b^{2}+c^{2})$,
\[
\begin{aligned}
\norm{g_t+e_t-e_{t+1}}^{2}
&\le 3\bigl(\norm{g_t}^{2}+\norm{e_t}^{2}+\norm{e_{t+1}}^{2}\bigr).
\end{aligned}
\tag{7}
\]

Taking expectation and applying variance decomposition,
\[
\EE[\norm{g_t}^{2}\mid w_t]
= \norm{\nabla f(w_t)}^{2} + \EE[\norm{g_t-\nabla f(w_t)}^{2}\mid w_t]
\le \norm{\nabla f(w_t)}^{2}+\sigma^{2}. \tag{8}
\]

Combine (7)–(8) and Lemma~\ref{lem:errorbound}:
\[
\EE\!\left[\norm{g_t+e_t-e_{t+1}}^{2}\mid w_t\right]
\le 3\Bigl(\norm{\nabla f(w_t)}^{2}+\sigma^{2}
+\tfrac{\sigma^{2}+2G^{2}}{\rho}\Bigr). \tag{9}
\]

\subsection*{Step 6: Assemble the inequality}
Insert (5)–(6) and (9) into (4):
\[
\begin{aligned}
\EE\!\left[f(w_{t+1})\mid w_t\right]
\le&
f(w_t)-\eta\norm{\nabla f(w_t)}^{2}
+2\eta G\sqrt{\tfrac{\sigma^{2}+2G^{2}}{\rho}}\\
&+\frac{3L\eta^{2}}{2}\norm{\nabla f(w_t)}^{2}
+\frac{3L\eta^{2}}{2}\Bigl(\sigma^{2}+\tfrac{\sigma^{2}+2G^{2}}{\rho}\Bigr).
\end{aligned}
\tag{10}
\]

\subsection*{Step 7: Choose stepsize to absorb the gradient term}
Require
\[
\eta\Bigl(1-\tfrac{3L\eta}{2}\Bigr)\ge\frac{\eta}{2},
\quad\text{i.e.}\quad
\eta\le\frac{1}{L\bigl(1+\sqrt{1-\rho}\bigr)}.
\tag{11}
\]
Under (11) we have $1-\frac{3L\eta}{2}\ge\frac{1}{2}$, so (10) simplifies to
\[
\EE\!\left[f(w_{t+1})\mid w_t\right]
\le f(w_t)-\frac{\eta}{2}\norm{\nabla f(w_t)}^{2}
+2\eta G\sqrt{\tfrac{\sigma^{2}+2G^{2}}{\rho}}
+\frac{3L\eta^{2}}{2}\Bigl(\sigma^{2}+\tfrac{\sigma^{2}+2G^{2}}{\rho}\Bigr). \tag{12}
\]

\subsection*{Step 8: Telescoping sum}
Take total expectation and sum (12) from $t=0$ to $T-1$:
\[
\begin{aligned}
\EE[f(w_T)]\le&
f(w_0)-\frac{\eta}{2}\sum_{t=0}^{T-1}\EE\!\left[\norm{\nabla f(w_t)}^{2}\right] \\
&+2\eta T G\sqrt{\tfrac{\sigma^{2}+2G^{2}}{\rho}}
+\frac{3L\eta^{2}T}{2}\Bigl(\sigma^{2}+\tfrac{\sigma^{2}+2G^{2}}{\rho}\Bigr).
\end{aligned}
\tag{13}
\]

Since $f(w_T)\ge f^{\star}$, rearrange (13):
\[
\frac{1}{T}\sum_{t=0}^{T-1}\EE\!\left[\norm{\nabla f(w_t)}^{2}\right]
\le
\frac{2\bigl(f(w_0)-f^{\star}\bigr)}{\eta T}
+4G\sqrt{\tfrac{\sigma^{2}+2G^{2}}{\rho}}
+\frac{3L\eta}{\;}\Bigl(\sigma^{2}+\tfrac{\sigma^{2}+2G^{2}}{\rho}\Bigr). \tag{14}
\]

\subsection*{Step 9: Simplify constants}
Observe that $4G\sqrt{(\sigma^{2}+2G^{2})/\rho}\le
\frac{2L\eta G^{2}}{\rho}+ \frac{2L\eta\sigma^{2}}{\rho}$ when $\eta\le 1/L$ (by $ab\le \tfrac{a^{2}}{2}+ \tfrac{b^{2}}{2}$). Applying this inequality to the second term of (14) yields
\[
\frac{1}{T}\sum_{t=0}^{T-1}\EE\!\left[\norm{\nabla f(w_t)}^{2}\right]
\le
\frac{2\bigl(f(w_0)-f^{\star}\bigr)}{\eta T}
+\frac{L\eta\sigma^{2}}{\rho}
+\frac{2L\eta G^{2}}{\rho}.
\tag{15}
\]

Equation (15) is exactly the bound stated in Theorem~\ref{thm:main}. ∎

\section*{Rate Derivation (With Constants)}
Choosing $\eta = \Theta\!\bigl(1/\sqrt{T}\bigr)$ (e.g. $\eta = \frac{c}{\sqrt{T}}$ with $c$ respecting (11)) gives
\[
\frac{1}{T}\sum_{t=0}^{T-1}\EE\!\left[\norm{\nabla f(w_t)}^{2}\right]
= O\!\left(\frac{1}{\sqrt{T}}\right)
+\underbrace{O\!\left(\frac{1}{\sqrt{T}}\right)}_{\text{compression term } \propto \tfrac{1}{\rho}}.
\]
Thus Top‑K EF‑SGD attains the optimal $O(1/\sqrt{T})$ rate for non‑convex stochastic optimization, with a multiplicative factor $1/\rho= d/K$ that quantifies the communication reduction.

\section*{Special Cases}
\begin{enumerate}[label=\alph*)]
    \item \textbf{Full communication} $K=d$ ($\rho=1$). The bound reduces to the classical SGD guarantee:
    \[
        \frac{1}{T}\sum_{t=0}^{T-1}\EE\!\left[\norm{\nabla f(w_t)}^{2}\right]
        \le\frac{2(f(w_0)-f^{\star})}{\eta T}+L\eta\sigma^{2}+2L\eta G^{2}.
    \]
    \item \textbf{Extreme compression} $K=1$ ($\rho=1/d$). The constants are inflated by a factor $d$, matching known lower bounds for unbiased compressors.
    \item \textbf{Diminishing stepsize} $\eta_t = \frac{c}{\sqrt{t+1}}$. Repeating the proof with a time‑varying $\eta_t$ yields the same $O(1/\sqrt{T})$ rate without requiring knowledge of $T$ in advance.
\end{enumerate}

\end{document}